\section{Evaluation Methodology}
\label{sec:methodology}

The evaluation addresses three questions: the formal-verifiability of inferred specifications under OpenJML's Extended Static Checker; the effectiveness of specification-guided test generation relative to baseline approaches; and the variation of inference effectiveness across method categories.

\subsection{Subject Selection}

Apache Commons Lang 3.14 was selected as the evaluation subject for its diversity of method types, extensive Javadoc documentation (which enables validation), maturity, deterministic pure functions (well suited to mutation testing), and precedent in the research literature~\cite{pacheco2007randoop, fraser2011evosuite}. Eleven classes covering five of the six method categories of Section~\ref{subsec:categorization} were selected: \texttt{NumberUtils}, \texttt{BooleanUtils}, \texttt{CharUtils} (utility); \texttt{MutableInt}, \texttt{MutableDouble}, \texttt{MutableBoolean} (state modification); \texttt{Pair}, \texttt{MutablePair} (factory/delegate); \texttt{Validate}, \texttt{Range} (control flow); and the \texttt{Mutable} interface (which provides the abstract operations realised by the three \texttt{Mutable*} classes; its declared methods are inherited and counted in the corresponding implementation classes). All public, package-private, and protected methods of these classes were included; no method was excluded a~posteriori. The \emph{other} category (callbacks and event handlers) is absent from the subset: Apache Commons Lang is a library of pure utility and value types and contains no methods of that kind. The 312 methods are distributed across the five remaining categories as shown in Table~\ref{tab:method-distribution}. The corresponding limitation for external validity is discussed in Section~\ref{sec:threats}.

\begin{table}[ht]
\centering
\caption{Distribution of methods by category in the Apache Commons Lang subset. Categories follow the taxonomy of Section~\ref{subsec:categorization}; the \emph{other} category contains zero methods because Apache Commons Lang does not include event-handler or callback methods.}
\label{tab:method-distribution}
\begin{tabular}{lrr}
\toprule
\textbf{Category} & \textbf{Count} & \textbf{\%} \\
\midrule
Utility Methods & 89 & 28.5\% \\
State Modification & 78 & 25.0\% \\
Accessors \& Mutators & 62 & 19.9\% \\
Control Flow & 48 & 15.4\% \\
Factory \& Delegate & 35 & 11.2\% \\
Other (callbacks, event handlers) & 0 & 0.0\% \\
\midrule
\textbf{Total} & \textbf{312} & \textbf{100\%} \\
\bottomrule
\end{tabular}
\end{table}

\subsection{Experimental Design}

The evaluation comprises four phases designed to isolate the contribution of formal specifications to test generation quality:

\begin{description}
    \item[Phase~1 (P1) --- Signature only:] test generation from method signatures alone.
    \item[Phase~2 (P2) --- Signature with guidance:] signatures supplemented with explicit guidance to cover edge cases.
    \item[Phase~3 (P3) --- Signature with specification:] signatures augmented with inferred formal specifications.
    \item[Phase~4 (P4) --- Source-code baseline:] full method source code, providing an upper bound on attainable quality.
\end{description}

Test generation used Google's Gemini~2.0 (Table~\ref{tab:llm-config}); a sampling temperature of 0.3 was selected to balance diversity against consistency. Five independent runs were performed per method per configuration to control for sampling non-determinism. Prompts were kept minimal and differed only in the inclusion of specifications or source code (Supplementary Material, Appendix~D).

\begin{table}[ht]
\centering
\caption{LLM configuration parameters}
\label{tab:llm-config}
\begin{tabular}{ll}
\toprule
\textbf{Parameter} & \textbf{Value} \\
\midrule
Model & Gemini 2.0 \\
Temperature & 0.3 \\
Max tokens & 4,096 \\
Top-p & 0.95 \\
Runs per configuration & 5 \\
\bottomrule
\end{tabular}
\end{table}

\subsection{Metrics}

\paragraph{Specification quality.}
The reported specification-quality metrics are the OpenJML \emph{discharge rate} (the proportion of inferred clauses formally verified by the Extended Static Checker against the implementation; Section~\ref{subsec:validation}) and the \emph{coverage rate} (the proportion of methods for which at least one non-trivial clause is emitted). Specifications are further classified by \emph{strength}: \emph{strong} when both precondition and postcondition are non-trivial, \emph{partial} when exactly one is non-trivial, and \emph{weak} when neither is. The strength classification is structural and is computed directly from the emitted JML, independently of the discharge outcome.

\paragraph{Test quality.} The reported metrics are test count, compilation rate, pass rate, and mutation score, the latter computed using PIT~\cite{coles2016pitest} (version 1.15.3) with the per-test timeout set to 10 seconds. The mutation operator set is the PIT \texttt{DEFAULTS} group, comprising \texttt{CONDITIONALS\_BOUNDARY} (replaces $<$ with $\leq$ and analogues), \texttt{INCREMENTS} (replaces \texttt{++} with \texttt{-{}-}), \texttt{INVERT\_NEGS} (negation of arithmetic expressions), \texttt{MATH} (binary arithmetic operator swaps), \texttt{NEGATE\_CONDITIONALS} (negation of boolean conditions), \texttt{RETURN\_VALS} (return-value replacement), \texttt{VOID\_METHOD\_CALLS} (removal of void calls), and \texttt{EMPTY\_RETURNS}, \texttt{FALSE\_RETURNS}, \texttt{TRUE\_RETURNS}, \texttt{NULL\_RETURNS}, \texttt{PRIMITIVE\_RETURNS} (return-value substitutions). Equivalent mutants are filtered automatically by PIT's default heuristics; the residual rate of trivially equivalent mutants is estimated at 4.7\% via the sampling procedure of Papadakis et al.~\cite{papadakis2019mutation} (200 random survived mutants per phase, dual-rater inspection, $\kappa = 0.83$). Mutation scores are reported on the post-filter mutant set.

The same 312 methods are used both for the OpenJML discharge analysis and for the test-generation effectiveness comparison reported in subsequent sections. The two analyses are methodologically independent observations on the same population: the discharge analysis is a verifier consistency check (is the spec formally consistent with the implementation?), whereas the effectiveness analysis compares LLM-generated tests across input conditions (does the spec, whatever its accuracy, change what the model produces?). No leakage between the two is possible: the LLM does not see the discharge outcome, and the discharge outcome is computed independently of the test-generation outcome.

\subsection{Statistical Analysis}

For each metric the analysis reports mean, standard deviation, median, interquartile range, and 95\% bootstrap confidence intervals (10{,}000 iterations). Phase comparisons use paired $t$-tests (or Wilcoxon signed-rank tests for non-normally distributed samples), with Cohen's $d$ as the effect-size measure. Bonferroni correction is applied across the full set of pairwise phase contrasts on the four reported metrics (test count, compilation rate, pass rate, mutation score): four phases yield $\binom{4}{2} = 6$ contrasts per metric, for $6 \times 4 = 24$ tests total at family-wise $\alpha = 0.05$ (per-test threshold $\alpha' = 0.05/24 \approx 0.0021$). Reported $p$-values that survive this corrected threshold are stated as $p < 0.001$ throughout.

\subsection{Comparison Baselines}

JML-Inferrer is compared against Jdoctor~\cite{jdoctor2018}, which uses natural-language processing to extract exception preconditions from Javadoc, and against direct LLM-based specification inference (the same Gemini~2.0 model prompted to infer JML from source code). The comparison reports coverage (the proportion of methods for which at least one non-trivial clause is emitted), OpenJML discharge rate where applicable, and run-to-run consistency. All comparisons use the same 312-method corpus, and all materials accompany the replication package.
